\documentclass[10pt,a4paper]{amsart}
\usepackage[margin=19mm]{geometry}
\usepackage{amsmath,amssymb,mathtools}
\usepackage[scr=rsfs,cal=euler]{mathalfa}
\usepackage{xfrac}
\usepackage[unicode,hidelinks,bookmarks=false]{hyperref}
\newcommand{\ie}{i.e.\ }
\newcommand{\cf}{cf.\ }
\newcommand{\resp}{respectively\ }
\newcommand{\Subsection}{\subsection}
\providecommand{\nobreakdash}{\nobreak\hbox{-}}
\setlength{\emergencystretch}{2em}
\multlinegap0pt
\usepackage{amssymb}
\newtheorem{corollary}{Corollary}
\newtheorem{theorem}{Theorem}
\title{Structural Redundancy in Subspace Network Coding via Atomic Decompositions}
\author{David Ramirez, Elvis Cabrera,, Jyrko Correa-Morris}
\date{Source: arXiv:2603.21390v1 (CC BY 4.0)}
\begin{document}
\maketitle

\noindent\emph{Source and licence.} Structural Redundancy in Subspace Network Coding via Atomic Decompositions. Version arXiv:2603.21390v1 (CC BY 4.0), distributed by arXiv under the Creative Commons Attribution 4.0 International licence, http://creativecommons.org/licenses/by/4.0/, as recorded in the arXiv licence metadata for this exact version and shown on the abstract page https://arxiv.org/abs/2603.21390v1. Full licence notice: \texttt{licenses/arxiv-2603.21390-CC-BY-4.0.txt}.\par\smallskip

\noindent\textit{Editorial scope.} Counted unit: the article's theorem thm:supermodular (source0.tex lines 424--436). The article's complete original proof of it is delivered with it (source0.tex lines 437--523, one \texttt{proof} environment of 87 lines). No mathematics is written, repaired or re-derived: the statement and the proof are the article's own lines, in the article's own notation, and no retained line is substituted or re-flowed.  Retained uncounted as the source-local context the proof needs: lines 289--290.  Nothing was omitted from any retained span, so the package carries no omission record.  No retained line contains a citation, so the wrapper carries no bibliography and no bibliography entry.  No retained line contains a figure environment, an image-inclusion command or drawing code, so no image is delivered and the package declares no asset dependency.  Editorial formatting only, in addition: an AMS \texttt{amsart} preamble that restates the article's own declarations and macros used by the retained lines, namely the environments \texttt{corollary}, \texttt{theorem} on separate counters, together with the standard packages those lines need.  The retained environments print in this document as Corollary 1, Theorem 1 in order of appearance, whereas the article prints the same items with the numbering of its own full document.  The complete native source of the article as distributed in the arXiv e-print for this exact version is bundled unchanged as \texttt{sources/196942/source0.tex}.
\begin{corollary}\label{Isotonic} $\mathfrak{N}$ is an isotone function; that is, for all subspaces $S,T\subseteq V$, $S\subseteq T$ implies $\mathfrak{N}(S)\le \mathfrak{N}(T)$.
\end{corollary}

\begin{theorem}\label{thm:supermodular}
For all subspaces $S,T \subseteq V$, the function $\mathfrak N$ satisfies
\begin{description}
     \item[(i)] $\mathfrak N(S) \cdot \mathfrak N(T)
\;\le\;
\mathfrak N(S \cap T) \cdot \mathfrak N(S + T)$, \quad (log-supermodularity)
\item[(ii)] $\mathfrak N(S) + \mathfrak N(T)
\;\le\;
\mathfrak N(S \cap T) + \mathfrak N(S + T).$ \quad (supermodularity)
\end{description}
    

\end{theorem}

\begin{proof}
    \item[(i)] We aim to prove that for all \( S, W \in L(V) \),
		\[
		\mathfrak{N}(S) \cdot \mathfrak{N}(W) \leq \mathfrak{N}(S + W) \cdot \mathfrak{N}(S \cap W),
		\]
		or equivalently,
		\[
		\frac{\mathfrak{N}(S + W)/\mathfrak{N}(W)}{\mathfrak{N}(S)/\mathfrak{N}(S \cap W)} \geq 1.
		\]
		
		Let us analyze both sides using the known formula:
		\[
		\mathfrak{N}(X) = \frac{\prod_{i=0}^{n_X - 1}(q^{n_X} - q^i)}{n_X! \cdot (q - 1)^{n_X}},
		\]
		where \( n_X = \dim(X) \). By factoring terms, we have:
		\begin{align*}
		\prod_{i=0}^{n_{S \vee W} - 1}(q^{n_{S \vee W}} - q^i)
		&= q^{n_W(n_{S \vee W} - n_W)} \cdot \prod_{i=0}^{n_W - 1}(q^{n_W} - q^i) \cdot \prod_{i=0}^{t - 1}(q^{n_{S \vee W}} - q^i), \\
		\prod_{i=0}^{n_S - 1}(q^{n_S} - q^i)
		&= q^{n_{S \wedge W}(n_S - n_{S \wedge W})} \cdot \prod_{i=0}^{n_{S \wedge W} - 1}(q^{n_{S \wedge W}} - q^i) \cdot \prod_{i=0}^{t - 1}(q^{n_S} - q^i),
		\end{align*}
		where \( t := n_{S \vee W} - n_W = n_S - n_{S \wedge W} \).
		
		Thus, the quotient becomes:
		\[
		\frac{\mathfrak{N}(S \vee W)/\mathfrak{N}(W)}{\mathfrak{N}(S)/\mathfrak{N}(S \wedge W)} = 
		\frac{
			q^{n_W t} \cdot n_W! \cdot \prod_{i=0}^{t - 1}(q^{n_{S \vee W}} - q^i)
		}{
			(q - 1)^t \cdot n_{S \vee W}!
		}
		\cdot
		\frac{
			(q - 1)^t \cdot n_S!
		}{
			q^{n_{S \wedge W} t} \cdot n_{S \wedge W}! \cdot \prod_{i=0}^{t - 1}(q^{n_S} - q^i)
		}.
		\]
		
		Now observe:
		\begin{itemize}
			\item The exponents satisfy \( n_W t - n_{S \wedge W} t = t(n_W - n_{S \wedge W}) \).
			\item Since \( n_{S \vee W} \geq n_S \), it follows that
			\[
			\prod_{i=0}^{t - 1}(q^{n_{S \vee W}} - q^i) \geq \prod_{i=0}^{t - 1}(q^{n_S} - q^i).
			\]
		\end{itemize}
		
		Therefore, the whole expression is bounded below by
		\[
		q^{t(n_W - n_{S \wedge W})} \cdot \frac{n_S! \cdot n_W!}{n_{S \vee W}! \cdot n_{S \wedge W}!}.
		\]
		
		Now we claim that for all \( t \leq n_S \),
		\[
		q^{t(n_W - n_S) + t^2} \geq \frac{(n_W + t)! \cdot (n_S - t)!}{n_W! \cdot n_S!}.
		\]
		We prove this by induction on \( t \).
		
		\emph{Base case:} For \( t = 0 \),
		\[
		q^0 = 1 = \frac{n_W! \cdot n_S!}{n_W! \cdot n_S!}.
		\]
		
		\emph{Inductive step:} Assume the inequality holds for \( t = k \):
		\[
		q^{k(n_W - n_S) + k^2} \geq \frac{(n_W + k)! \cdot (n_S - k)!}{n_W! \cdot n_S!}.
		\]
		
		We must prove it for \( t = k + 1 \), that is:
		\[
		q^{(k+1)(n_W - n_S) + (k+1)^2} \geq \frac{(n_W + k + 1)! \cdot (n_S - k - 1)!}{n_W! \cdot n_S!}.
		\]
		
		Dividing both sides of the \( t = k+1 \) inequality by the inductive hypothesis, it suffices to show:
		\[
		q^{(n_W + k + 1) - (n_S - k)} \geq \frac{n_W + k + 1}{n_S - k}.
		\]
		
		Since \( q \geq 2 \), the inequality
		\[
		q^{a-b} \geq \frac{a}{b}
		\]
		holds for all \( a,b \in \mathbb{N} \) such that $a>b$, and thus the desired inequality follows. 

        \item[(ii)] It is a well-known fact that if a function is isotone Corollary \ref{Isotonic} and log-supermodular, then it is also supermodular.
\end{proof}

\end{document}
